\documentclass[11pt]{article}
\usepackage[paperwidth=210mm,paperheight=297mm,margin=16mm]{geometry}
\usepackage{fontspec,amsmath,array,multirow,graphicx}
\usepackage{polyglossia}
\setmainlanguage{latin}\setotherlanguages{arabic,greek}
\setmainfont{Linux Libertine O}[Ligatures=TeX]
\newfontfamily\arabicfont{Amiri}[Script=Arabic]
\newfontfamily\greekfont{FreeSerif}[Script=Greek]
\newfontfamily\NallinoSigns{NallinoSigns.otf}[Path=./fonts/]
\newcommand{\AbjadZero}{{\NallinoSigns\symbol{"E000}}}
\newcommand{\AbjadOver}[1]{\leavevmode\vbox{\hrule height .45pt\kern1.2pt\hbox{#1}}}
\newcommand{\Name}[1]{{\addfontfeatures{LetterSpace=12}#1}}
\pagestyle{empty}\setlength{\parindent}{0pt}
\renewcommand{\arraystretch}{1.18}

\begin{document}
\clearpage
{\large \textarabic{٢٢٨}}\par\vspace{1mm}
\begin{Arabic}
\begin{center}بسم\textsuperscript{1} الله الرحمن الرحيم صلّى الله على سيّدنا محمّد رسوله الكريم وعلى آله وصحبه وسلّم\end{center}\vspace{-2mm}
\noindent\fbox{\parbox{\dimexpr\textwidth-2\fboxsep-2\fboxrule}{\centering\large جَدْوَل تأْريخ الملوك اليونانيّة من لدن بُخْتَنَصَّرَ الاوّل ومنه بتأْريخ المِجِسْطِي}}\par\vspace{0.5mm}
\noindent\begin{minipage}[t]{0.495\textwidth}\setlength{\tabcolsep}{2.5pt}\begin{tabular}{|p{44mm}|>{\centering\arraybackslash}p{7mm}|>{\centering\arraybackslash}p{9mm}|}\hline
\centering اسماء الملوك & \rotatebox{90}{\small عدد ما ملكوا} & \rotatebox{90}{\small مجموع السنين} \\ \hline
{}بُخْتَنَصَّر الاوّل & يد & يد \\
{}نديوس\textsuperscript{2} & ب & يو \\
{}خنزىروس\textsuperscript{3} & ه & كا \\
{}الميوا\textsuperscript{4} & ه & كو \\
{}مردوقنفذ\textsuperscript{5} & يب & لح \\
{}ارقينوا\textsuperscript{6} & ه & مج \\
{}افسيليطوس الاوّل & ب & مه \\
{}بيل\textsuperscript{7} & ج & مح \\
{}افرانديوس\textsuperscript{8} & و & ند \\
{}ارسفل\textsuperscript{9} & ا & نه \\
{}افسيلطوس الثاني & د & نط \\
{}مسيسيموردقس\textsuperscript{10} & ح & صز \\
{}اردوسون\textsuperscript{11} & يج & ف \\
{}سسدوخنس\textsuperscript{12} & ك & ق \\
{}قنيلدنوس\textsuperscript{13} & كب & قكب \\
{}بُخْتَنَصَّر الثاني & كا & قمج \\
\hline\end{tabular}\end{minipage}\hfill\begin{minipage}[t]{0.495\textwidth}\setlength{\tabcolsep}{2.5pt}\begin{tabular}{|p{44mm}|>{\centering\arraybackslash}p{7mm}|>{\centering\arraybackslash}p{9mm}|}\hline
\centering اسماء الملوك & \rotatebox{90}{\small عدد ما ملكوا} & \rotatebox{90}{\small مجموع السنين} \\ \hline
{}بُخْتَنَصَّر الثالث وهو فاتح المَقْدِس & مج & قفو \\
{}برعليا ليقوا\textsuperscript{14} & ب & قفح \\
{}بلشصَّر & د & قضب \\
{}دريوس الآذَرِيّ & يز & رط \\
{}كورش & ط & ريح \\
{}قمبوسوس & ح & ركو \\
{}دريوس & لو & رب \\
{}اخشقيوس\textsuperscript{15} & كا & رفج \\
{}ارطَحْشَسْت\textsuperscript{16} الاوّل & مج & سكو \\
{}دريوس\textsuperscript{17} & يط & سمه \\
{}ارطَحْشَسْت\textsuperscript{18} الثاني & مو & سضا \\
{}اخوس\textsuperscript{19} & كا & تيب \\
{}غيرون\textsuperscript{20} & ى & تكب \\
{}دريوس ارسخ & و & تكح \\
{}الاسكندر الماقذونيّ & و & تلد \\
\hline\end{tabular}\end{minipage}
\end{Arabic}
\par\vspace{3mm}\noindent\rule{40mm}{.4pt}\par\vspace{1mm}{\small 1) Incipit f. 154,v. --- 2) Cod. \textarabic{نديوب} --- 3) Cod. \textarabic{حيرنقون}; apud \Name{al-Bīrūnī}, \textit{Chron.} \textarabic{٨٨}, ubi eadem regum Babyloniensium tabula ex Canone Ptolemaei legitur, \textarabic{حرىفون} --- 4) Sic; \Name{al-Bīrūnī} \textarabic{اىلوعىو} --- 5) Cod. \textarabic{مودوفيقد} --- 6) Cod. \textarabic{ارقنيوا}; \Name{al-Bīrūnī} \textarabic{اريقينو} --- 7) \Name{Al-Bīrūnī} \textarabic{بيل بيس} --- 8) Cod. \textarabic{افراندولن}; \Name{al-Bīrūnī} \textarabic{اوڡراىدىدر} --- 9) Sic; \Name{al-Bīrūnī} \textarabic{ارسعل} --- 10) Cod. \textarabic{ميسليموردقس}; \Name{al-Bīr.} \textarabic{سسلموردقش} --- 11) Sic; \Name{al-Bīr.} \textarabic{ارديدينو} --- 12) Cod. \textarabic{سسدوحس}; \Name{al-Bīr.} \textarabic{سسدوكن} --- 13) Cod. \textarabic{فنيلدوس}; \Name{al-Bīr.} \textarabic{ڡلسرورڡلدن} --- 14) Sic; \Name{al-Bīr.} \textarabic{برخلالتغر} --- 15) Sic; \Name{al-Bīr.} \textarabic{اخشيرش} --- 16) Cod. sine punctis. --- 17) Cod. \textarabic{طريوش} --- 18) Cod. sine punctis. --- 19) Cod. \textarabic{اموس} --- 20) Sic; \Name{al-Bīr.} \textarabic{وىرون}}
\clearpage
{\large \textarabic{٢٢٩}}\par\vspace{1mm}
\begin{Arabic}
\noindent\begin{minipage}[t]{0.495\textwidth}\setlength{\tabcolsep}{2.5pt}\begin{tabular}{|p{44mm}|>{\centering\arraybackslash}p{7mm}|>{\centering\arraybackslash}p{9mm}|}\hline
\centering اسماء الملوك & \rotatebox{90}{\small عدد ما ملكوا} & \rotatebox{90}{\small مجموع السنين} \\ \hline
\hline\multicolumn{3}{|c|}{وبعد هذا ملوك اتور\textsuperscript{1}} \\ \hline
{}فيافوس\textsuperscript{2} ابو ذي القَرْنَيْن & ز & ز \\
{}الاسكندر البَنَّاء & يب & يط \\
{}بطلميوس أَرْنَب\textsuperscript{3} & ك & لط \\
{}بطلميوس فيلاذلفوس & لح & عز \\
{}بطلميوس اورغطس الاوّل & كه & قب \\
{}بطلميوس فيليفطور & يز & قيط \\
{}بطلميوس افيفنيس & كد & قمج \\
{}بطلميوس فيليميطور & له & قعح \\
{}بطلميوس اورغطس الثاني & كط & رز \\
{}بطلميوس سوطر & لو & رمج \\
{}بطلميوس\textsuperscript{4} ديونسيوس & كط & رعب \\
{}بطلميوس قلوبطرا & كب & رضد \\
\hline\multicolumn{3}{|c|}{وبعد هذا ملوك الروم} \\ \hline
{}اغسطس & مج & سلز \\
{}طيبريوس & كب & سنط \\
{}غائيوس\textsuperscript{5} & د & سصج \\
{}قلوذيوس & يد & سعز \\
\hline\end{tabular}\end{minipage}\hfill\begin{minipage}[t]{0.495\textwidth}\setlength{\tabcolsep}{2.5pt}\begin{tabular}{|p{44mm}|>{\centering\arraybackslash}p{7mm}|>{\centering\arraybackslash}p{9mm}|}\hline
\centering اسماء الملوك & \rotatebox{90}{\small عدد ما ملكوا} & \rotatebox{90}{\small مجموع السنين} \\ \hline
{}نارون & يد & سضا \\
{}ويسفسيانوس\textsuperscript{6} & ى & تا \\
{}طيطوس\textsuperscript{7} & ج & تد \\
{}دومطيانوس\textsuperscript{8} & يه & تيط \\
{}ناروس\textsuperscript{9} & ا & تك \\
{}طرايانوس\textsuperscript{10} & يط & تلط \\
{}اذريانس & كا & تص \\
{}انطونينس\textsuperscript{11} & كج & تفج \\
{}قومذوس & لب & ثيه \\
{}ساوىرس\textsuperscript{12} & كه & ثم \\
{}انطونينوس وَحْدَه\textsuperscript{13} & د & ثمد \\
{}الكسندرس برغما\textsuperscript{14} & يج & ثنز \\
{}مكسميانوس\textsuperscript{15} & ج & ثص \\
{}غورديانوس & و & ثصو \\
{}فيلفس & و & ثعب \\
{}ذاقيانس & ا & ثعج \\
{}غالس & ج & ثعو \\
{}والرينس\textsuperscript{16} & يه & ثضا \\
\hline\end{tabular}\end{minipage}
\end{Arabic}
\par\vspace{3mm}\noindent\rule{40mm}{.4pt}\par\vspace{1mm}{\small 1) Error amanuensis; in titulo f. 155,r. recte legitur \textarabic{الملوك اليونانية} --- 2) Cod. \textarabic{قيلقوس} --- 3) Cod. \textarabic{اربيا}. Cognomina sequentium Ptolemaeorum hoc modo in cod. deformata leguntur: \textarabic{قيذايقش}, \textarabic{اوعسطس}, \textarabic{قيلنقطوس}, \textarabic{افيقينس}, \textarabic{قنليقطور}, \textarabic{اوعسطوس}, \textarabic{سرسوطر}, \textarabic{ديولسلير} --- 4) Incipit f. 155,r., cui titulus \textarabic{جدول تاريخ الملوك اليونانية والرومية من لدن قيلقوس} --- 5) Cod. \textarabic{غالبوس}; cf. \Name{al-Bīr.} \textarabic{٩٢} et \Name{al-Mas‘ūdī}, \textit{Tanbīh} \textarabic{١٢٢}, adn. --- 6) Cod. \textarabic{ويسقيسرموس} --- 7) Cod. \textarabic{طنطوس} --- 8) Cod. \textarabic{ذومطيناوس} --- 9) Cod. \textarabic{ثادوس} --- 10) Cod. \textarabic{طوايانوس} --- 11) Cod. \textarabic{انطوينيش} --- 12) Cod. \textarabic{ساوطيلس} --- 13) Cod. \textarabic{انطوبلبوس وجده} --- 14) Sic; pro \textarabic{بن ماميا}. Codd. \Name{al-Bīr.} \textarabic{٩٤} \textarabic{بزيما}; \Name{Al-Mas‘ūdī}, \textit{Prairies} II, 306 \textarabic{الاسكندر مامياس}; \textit{Tanbīh} \textarabic{١٣٣} \textarabic{الاسكندرس يلقَّب مامياس}; \Name{Abū ’l-farag}, ed. Ṣālḥānī \textarabic{١٢٦} \textarabic{الكسندروس بن ماما} etc. --- 15) Cod. \textarabic{مكسابوس} --- 16) Cod. \textarabic{وارثيس}}
\clearpage
{\large \textarabic{٢٣٠}}\par\vspace{1mm}
\begin{Arabic}
\noindent\begin{minipage}[t]{0.495\textwidth}\setlength{\tabcolsep}{2.5pt}\begin{tabular}{|p{44mm}|>{\centering\arraybackslash}p{7mm}|>{\centering\arraybackslash}p{9mm}|}\hline
\centering اسماء الملوك & \rotatebox{90}{\small عدد ما ملكوا} & \rotatebox{90}{\small مجموع السنين} \\ \hline
{}قلوذيوس & ا & ثضب \\
{}اوريلينوس\textsuperscript{1} & و & ثضح \\
{}فربوس\textsuperscript{2} & ز & خه \\
{}قاروس وقارينوس\textsuperscript{3} & ب & خز \\
{}دقلطيانوس & كا & خكح \\
\hline\multicolumn{3}{|c|}{وبعد هذا ملوك النَّصْرانيّة} \\ \hline
{}قوسطنطينس & لب & خص \\
{}قوسطوس & كد & خفد \\
{}يوليانس الحَنِيف & ب & خفو \\
{}يوبيانس\textsuperscript{4} & ا & خفز \\
{}ثاذوسيوس\textsuperscript{5} & يد & ذا \\
{}وليس\textsuperscript{6} & يز & ذيح \\
{}ارقاذيوس & يج & ذلا \\
{}ثاوذسيوس\textsuperscript{7} & مب & ذعج \\
{}مارقيانس & و & ذعط \\
{}لأون & يح & ذضز \\
{}زينون\textsuperscript{8} & يز & ضد \\
\hline\end{tabular}\end{minipage}\hfill\begin{minipage}[t]{0.495\textwidth}\setlength{\tabcolsep}{2.5pt}\begin{tabular}{|p{44mm}|>{\centering\arraybackslash}p{7mm}|>{\centering\arraybackslash}p{9mm}|}\hline
\centering اسماء الملوك & \rotatebox{90}{\small عدد ما ملكوا} & \rotatebox{90}{\small مجموع السنين} \\ \hline
{}انسطاسيوس\textsuperscript{9} & كز & ضما \\
{}يوسطينس الاوّل & ط & ضن \\
{}يوسطينينوس\textsuperscript{10} & لز & ضفز \\
{}يوسطينس الثاني & يد & ظا \\
{}طيبريوس\textsuperscript{11} & د & ظه \\
{}ماورقيوس\textsuperscript{12} & ك & ظكه \\
{}فوقاس & ح & ظلج \\
{}هرقلس صاحب العرب & لا & ظصد \\
{}قسطنطنس & ا & ظصه \\
{}قوسطنطيوس\textsuperscript{13} & كز & ظضب \\
{}لسطبوس\textsuperscript{14} & يو & ا ح \\
{}طيبريوس\textsuperscript{15} الثاني & ى & ا يح \\
{}يوسطينس & ج & ا كا \\
{}ڡيقوس\textsuperscript{16} & ز & ا كح \\
{}يوسطينينس\textsuperscript{17} & و & ا لد \\
{}فيلفقوس\textsuperscript{18} & ج & ا لز \\
{}انسطاس\textsuperscript{19} & ب & ا لط \\
{}ثاوذسيوس\textsuperscript{20} & ا & ا م \\
\hline\end{tabular}\end{minipage}
\end{Arabic}
\par\vspace{3mm}\noindent\rule{40mm}{.4pt}\par\vspace{1mm}{\small 1) Cod. \textarabic{اقوتلينوس} --- 2) Cod. \textarabic{قربوس} --- 3) Cod. \textarabic{قارسوس} --- 4) Cod. \textarabic{بوسانس} --- 5) Cod. \textarabic{نادرسوس} --- 6) Cod. \textarabic{ووليس} --- 7) Cod. \textarabic{ناوذسيوس} --- 8) Incipit f. 155,v., cui titulus: \textarabic{بقية جداول تاريخ ملوك النصرانية}. In hac pag. amanuensis oblitus est litteras numerales orientales archetypi in maghrebinicas convertere; ergo hic \textarabic{ض} = 800, \textarabic{ظ} = 900. --- 9) Cod. \textarabic{اسطاثبوس} --- 10) Cod. \textarabic{بوقݐطݐوس} --- 11) Cod. \textarabic{طنيربوس} --- 12) Cod. \textarabic{ناورنقوس} --- 13) Cod. \textarabic{قوسطنطينوس} --- 14) Sic. --- 15) Cod. \textarabic{طنيربوس} --- 16) Sic. --- 17) Cod. \textarabic{قوسطنس} --- 18) Cod. \textarabic{قيقوس} --- 19) Cod. \textarabic{السطاس} --- 20) Cod. \textarabic{تاودسبونس}}
\clearpage
{\large \textarabic{٢٣١}}\par\vspace{1mm}
\begin{Arabic}
\noindent\fbox{\parbox{\dimexpr\textwidth-2\fboxsep-2\fboxrule}{\centering\large جَدْوَل ما بين التواريخ}}\par\vspace{0.5mm}
\noindent\begin{tabular}{|p{\dimexpr\textwidth-22mm\relax}|>{\centering\arraybackslash}p{12mm}|}\hline
{}بين مُلْك بُخْتَنَصَّر الاوّل الى ممات الاسكندر الماقذونيّ من السنين المصريّة & تكد \\
{}ثمّ مَلَكَ بعد ذلك فيلبوس ابو ذي القرنَيْن فمن ملكه الى ملك اغسطس الروميّ & رضد \\
{}ومن ملك اغسطس الروميّ الى ملك دقلطيانوس وهو من ملوك النصرانيّة & سيج \\
{}ومن ملك دقلطيانوس الى ملك يليانس الحنيف & عز \\
{}ثمّ ملك يليانس وعاد الملك الى النصرانيّة وثمّ يملك بلياس سور\textsuperscript{1} & ب \\
{}ومن ملك دقلطيانوس الى ملك هرقل صاحب العرب & سكو \\
{}ثمّ ملكت العرب فإنّ من هجرة النبيّ الى دولة معاوية الأُمَوِيّ & لط \\
{}والى أن صار الامر لبني العبّاس وبَقِيَ في بني العبّاس & قكز \\
\hline\end{tabular}
\end{Arabic}
\par\vspace{5mm}
\begin{Arabic}
\noindent\fbox{\parbox{\dimexpr\textwidth-2\fboxsep-2\fboxrule}{\centering\large جَدْوَل\textsuperscript{2} تأْريخ الخُلَفَاء من لَدُنِ الهِجْرة هجرة النبيّ صلَّى الله عليه وسلّم}}\par\vspace{0.5mm}
\setlength{\tabcolsep}{2pt}\noindent\begin{tabular}{|b{\dimexpr\textwidth-60mm\relax}|>{\centering\arraybackslash}p{6.5mm}|>{\centering\arraybackslash}p{6.5mm}|>{\centering\arraybackslash}p{6.5mm}||>{\centering\arraybackslash}p{6.5mm}|>{\centering\arraybackslash}p{6.5mm}|>{\centering\arraybackslash}p{6.5mm}|}\hline
\parbox[b]{\dimexpr\textwidth-60mm\relax}{\centering\small أسْماء الخُلَفَاء الراشدين من لدن الهجرة على أنّ اوّل يوم من المحرَّم\newline الجُمْعة والذي يُعْمَل عليه في التأْريخ الخميس وهذا المحرَّم لاوّل\newline سنة الهجرة} & \multicolumn{3}{c||}{\small ما ملك كلّ واحد منهم} & \multicolumn{3}{c|}{\small مجموعة السنين} \\ \cline{2-7}
 & {\scriptsize سنون} & {\scriptsize شهور} & {\scriptsize ايّام} & {\scriptsize سنون} & {\scriptsize شهور} & {\scriptsize ايّام} \\ \hline
{}كانت هِجْرة النبيّ محمَّد صلَّى الله عليه وسلّم من مَكَّة\newline الى المدينة سنة احدى لها & \AbjadZero{} & ب & ح & \AbjadZero{} & ب & ح \\
{}فمكَث مُهاجِرًا بالمدينة حتَّى قُبِضَ & ط & يا & كب & ى & ب & \AbjadZero{} \\
{}ابو بَكْر بن ابي قُحافَةَ من بني تَيْم & ب & ج & ح & يب & د & ح \\
{}عُمَر بن الخَطّاب من بني عَدِيّ & ى & و & يز & كب & يا & كه \\
{}وكانت الشَّوْرَى بعد عمر بن الخطّاب & \AbjadZero{} & \AbjadZero{} & ج & كب & يا & كح \\
{}عُثْمان بن عَفّانَ من بني أُمَيَّةَ & يا & يا & يط & لد & يا & يز \\
{}عَلِيّ بن ابي طالِب والفتْنة & د & ط & \AbjadZero{} & لط & ح & يو \\
\hline\end{tabular}
\end{Arabic}
\par\vspace{3mm}\noindent\rule{40mm}{.4pt}\par\vspace{1mm}{\small 1) Sic. --- 2) Inc. f. 156,r.}
\clearpage
{\large \textarabic{٢٣٢}}\par\vspace{1mm}
\begin{Arabic}
\setlength{\tabcolsep}{2pt}\noindent\begin{tabular}{|b{\dimexpr\textwidth-60mm\relax}|>{\centering\arraybackslash}p{6.5mm}|>{\centering\arraybackslash}p{6.5mm}|>{\centering\arraybackslash}p{6.5mm}||>{\centering\arraybackslash}p{6.5mm}|>{\centering\arraybackslash}p{6.5mm}|>{\centering\arraybackslash}p{6.5mm}|}\hline
\parbox[b]{\dimexpr\textwidth-60mm\relax}{\centering\small أسْماء الخُلَفَاء الراشدين\textsuperscript{1} من لدن الهجرة} & \multicolumn{3}{c||}{\small ما ملك كلّ واحد منهم} & \multicolumn{3}{c|}{\small مجموعة السنين} \\ \cline{2-7}
 & {\scriptsize سنون} & {\scriptsize شهور} & {\scriptsize ايّام} & {\scriptsize سنون} & {\scriptsize شهور} & {\scriptsize ايّام} \\ \hline
{}وإلى بَيْعة مُعاوِيَةَ بن ابي سُفْيانَ & \AbjadZero{} & و & ج & م & ب & ك \\
{}معاوية بن ابي سفيان بن حَرْب بن أُمَيَّةَ & يط & ج & كه & نط & و & يه \\
{}يَزيد بن معاوية بن ابي سُفْيان & ج & ح & \AbjadZero{} & صج & ب & يه \\
{}معاوية بن يزيد بن معاوية & \AbjadZero{} & ج & كب & صج & و & ز \\
{}عبد الله بن الزُّبَيْر ومَرْوان بن الحَكَم & \AbjadZero{} & د & \AbjadZero{} & صج & ى & ز \\
{}عبد الله بن الزبير من بني أَسَدِ & ح & ه & \AbjadZero{} & عب & ج & ز \\
{}عبد الملك بن مَرْوانَ حتّى قُتِلَ ابن الزُّبَيْر & ا & ب & ج & عج & ه & ى \\
{}عبد الملك بن مروان بن الحَكَم & يب & د & و & فه & ط & يه \\
{}الوَلِيد بن عبد الملك بن مَرْوانَ & ط & ز & كط & ضه & ه & يد \\
{}سُلَيْمان بن عبد الملك & ب & ز & كط & ضح & ا & يج \\
{}عمر بن عبد العزيز بن مروان & ب & ه & يج & ق & و & كو \\
{}يَزيد\textsuperscript{2} بن عبد الملك بن مروان & د & \AbjadZero{} & ا & قد & و & كد \\
{}هِشام بن عبد الملك بن مروان & يط & ح & ط & قكد & ج & و \\
{}الوليد بن يَزيد بن عبد الملك بن مروان & ا & ب & كا & قكه & ه & كز \\
{}وكانت الفِتْنة بعد قَتْل الوليد & \AbjadZero{} & ب & كه & قكه & ح & كب \\
{}يزيد بن الوليد بن عبد الملك & \AbjadZero{} & ب & ط & قكه & يا & ا \\
{}ابراهيم بن الوليد بن عبد المَلك & \AbjadZero{} & ب & يا & قكو & ا & يب \\
{}مروان بن محمَّد بن مروان حتَّى قُتِلَ & ه & ب & \AbjadZero{} & قلا & ج & يب \\
{}ثمّ عاد الأَمْر لبنى هاشم &  &  &  &  &  &  \\
{}ابو العَبَّاس عبد الله بن محمَّد السَّفَّاح & د & ح & ب & قله & يا & يد \\
{}وحتَّى انتهت البَيْعة الى ابي جَعْفَر & \AbjadZero{} & \AbjadZero{} & يد & قله & يا & كح \\
\hline\end{tabular}
\end{Arabic}
\par\vspace{3mm}\noindent\rule{40mm}{.4pt}\par\vspace{1mm}{\small 1) Hic et in pag. seq. \textarabic{الراشدين} quod cod. semper habet retineo. --- 2) Inc. f. 156,v.}
\clearpage
{\large \textarabic{٢٣٣}}\par\vspace{1mm}
\begin{Arabic}
\setlength{\tabcolsep}{2pt}\noindent\begin{tabular}{|b{\dimexpr\textwidth-60mm\relax}|>{\centering\arraybackslash}p{6.5mm}|>{\centering\arraybackslash}p{6.5mm}|>{\centering\arraybackslash}p{6.5mm}||>{\centering\arraybackslash}p{6.5mm}|>{\centering\arraybackslash}p{6.5mm}|>{\centering\arraybackslash}p{6.5mm}|}\hline
\parbox[b]{\dimexpr\textwidth-60mm\relax}{\centering\small أسْماء الخُلَفَاء الراشدين من لدن الهجرة} & \multicolumn{3}{c||}{\small ما ملك كلّ واحد منهم} & \multicolumn{3}{c|}{\small مجموعة السنين} \\ \cline{2-7}
 & {\scriptsize سنون} & {\scriptsize شهور} & {\scriptsize ايّام} & {\scriptsize سنون} & {\scriptsize شهور} & {\scriptsize ايّام} \\ \hline
{}ابو جَعْفَرِ المنصور عبد الله بن محمَّد & كا & يا & ح & قنز & يا & و \\
{}وحتّى انتهى الخَبَر الى المَهْدِيّ & \AbjadZero{} & \AbjadZero{} & يب & قنز & يا & يح \\
{}المَهْدِيّ محمَّد بن ابي جَعْفَرِ المنصور & ى & ا & ه & قصح & \AbjadZero{} & يج \\
{}وحتّى انتهى الخَبَر الى موسى بن المهديّ & \AbjadZero{} & \AbjadZero{} & ح & قصح & ا & ا \\
{}الصادق مُوسَى بن محمَّد المَهْدِيّ & ا & ا & يه & قصط & ب & يو \\
{}الرَّشِيد هارون بن محمَّد المهديّ & كج & ب & يو & قصب & ه & ج \\
{}وحتّى انتهى الخَبَر الى محمَّد بن هارونَ & \AbjadZero{} & \AbjadZero{} & يب & قصب & ه & يه \\
{}الأمِين محمَّد بن الرشيد حتَّى خُلِعَ وحُبِسَ & ج & \AbjadZero{} & كه & قصه & و & ى \\
{}فمكَث محبوسًا & \AbjadZero{} & \AbjadZero{} & ب & قصه & و & يب \\
{}ثمّ أُخْرِجَ وبُويِعَ وحارب وحُوصِرَ حتَّى قُتِلَ & ا & و & يج & قصز & \AbjadZero{} & كه \\
{}المَأْمون عبد الله بن هارونَ الرَّشِيد & ك & ه & كب & ريز & و & يز \\
{}المُعْتَصِم محمَّد بن هارون الرشيد & ح & ه & ب & ركو & ب & يط \\
{}الواثِق\textsuperscript{1} بالله هارون بن محمَّد المعتصِم & ه & ط & ه & رلا & يا & كد \\
{}المُتَوَكِّل على الله جَعْفَر بن محمَّد المعتصِم & يد & ط & ط & رمو & ط & ج \\
{}المُنْتَصِر بالله محمَّد بن المتوكِّل & \AbjadZero{} & و & \AbjadZero{} & رمز & ج & ج \\
{}المُسْتَعِين بالله الى أن انحدر الى مدينة السَّلام & ب & ط & ج & رن & \AbjadZero{} & و \\
{}والى أن بُويِعَ المُعْتَزّ بالله بسُرَّ مَنْ رَأَى\textsuperscript{2} & \AbjadZero{} & \AbjadZero{} & ح & رن & \AbjadZero{} & يد \\
{}والى ان خُطِبَ المعتزّ بالله بمدينة السلام & \AbjadZero{} & يا & ك & رنا & \AbjadZero{} & د \\
{}والى ان خُلِعَ المعتزّ بالله & ج & و & كج & رند & و & كز \\
{}والى ان بُويِعَ المُهْتَدِي بالله & \AbjadZero{} & \AbjadZero{} & ب & رند & و & كط \\
{}المُهْتَدِي بالله بن الواثق بالله & \AbjadZero{} & يا & يح & رنه & و & يز \\
\hline\end{tabular}
\end{Arabic}
\par\vspace{3mm}\noindent\rule{40mm}{.4pt}\par\vspace{1mm}{\small 1) Inc. f. 157,r., cui titulus: \textarabic{جدول تاريخ الخلفاء من لدن هجرة النبي صلى الله عليه وسلم} --- 2) Cod. \textarabic{بشر بن رانى}}
\par\vfill\hfill 30
\clearpage
{\large \textarabic{٢٣٤}}\par\vspace{1mm}
\begin{Arabic}
\setlength{\tabcolsep}{2pt}\noindent\begin{tabular}{|b{\dimexpr\textwidth-60mm\relax}|>{\centering\arraybackslash}p{6.5mm}|>{\centering\arraybackslash}p{6.5mm}|>{\centering\arraybackslash}p{6.5mm}||>{\centering\arraybackslash}p{6.5mm}|>{\centering\arraybackslash}p{6.5mm}|>{\centering\arraybackslash}p{6.5mm}|}\hline
\parbox[b]{\dimexpr\textwidth-60mm\relax}{\centering\small أسْماء الخُلَفَاء الراشدين من لدن الهجرة} & \multicolumn{3}{c||}{\small ما ملك كلّ واحد منهم} & \multicolumn{3}{c|}{\small مجموعة السنين} \\ \cline{2-7}
 & {\scriptsize سنون} & {\scriptsize شهور} & {\scriptsize ايّام} & {\scriptsize سنون} & {\scriptsize شهور} & {\scriptsize ايّام} \\ \hline
{}المُعْتَمِد على الله احمد بن المُتَوَكِّل & كح & \AbjadZero{} & ج & رعح & و & ك \\
{}المُعْتَضِد بالله احمد بن المُوَفَّق & ط & ط & ب & رفح & ج & كب \\
{}المُكْتَفِي بالله عليّ بن المُعْتَضِد & ز & و & د & رضد & ط & كو \\
{}المُقْتَدِر بالله جَعْفَر بن أَحْمَد & كد & يا & يه & سيط & ج & يح \\
{}القاهِر بالله محمَّد بن احمد & ا & و & و & سك & ى & كد \\
{}الرَّاضِي بالله محمَّد بن جَعْفَرِ & و & ى & ط & سكز & ط & ج \\
{}المُتَّقِي لله ابراهيم بن جعفر & ج & يا & د & سلا & ح & ز \\
{}المُسْتَكْفِي بالله عبد الله بن عليّ & ا & و & يه & سلج & ب & كب \\
{}المُطِيع لله الفَضْل بن جَعْفَرِ &  &  &  &  &  &  \\
\hline\end{tabular}
\end{Arabic}
\end{document}
